\section{讨论与限制}

下面七条加本篇各一条，构成本篇的全部已知边界。其中 L4 把团队要求的
「幂律主张在多大样本量上必须放弃」写成可复算的判据，L1 是本研究的首要威胁，
其余各条说明这两条为什么不是工程瑕疵。

\begin{itemize}
\item \textbf{L1 同源自证，这是首要威胁，且与第 6 节的机制断言直接冲突。}
事件类别标签与人物名出自同一部原书：被统计的叙事与被统计的分类可能是同一段文字的两种编码。
本篇的处理是：只报分布形状与时期占比的位移，把位移叫相对基线的差，不叫效应量；
凡是把这里的形状当成明代社会真实结构的读法，本篇一律不做。
\item \textbf{L2 关系边不等于社会关系。} 关系表里一条「敌对冲突」是叙述者写下的关系，
不是双方真的敌对过；它可能是同一件事、同一句话、同一场战役被写成一条边。
本篇只量边被标成了什么类别，不量类别背后发生过什么。
\item \textbf{L3 本篇刻意不做词频与热度那一层。}
本篇没有做原书全文的词频统计，也没有用出场章节数当"热度中心性"的代理；
原因是同类研究已经证明词频与热度都可能只是行文与编排习惯，
它们不能证明结构。本篇因此把可断言的范围钉在四件可点数的事上：
基尼、连通块、事件年份分布、类别占比迁移。任何"时人感到崩塌"
"叙事者预感亡国"一类的断言，都不在本篇证据之内。
\item \textbf{L4 幂律主张的样本量上限（本篇的硬边界）。}
按 \citet{clauset2009} 的口径，幂律主张至少要求两件事同时成立：
尾指数在多个 $x_{\min}$ 上稳定、尾样本足够大。本篇 $x_{\min}=10$ 时尾样本
只有 142 条，$x_{\min}=15$ 时只剩 98 条。经验地讲：
在这批数据上，只要尾样本量低于约一百五十条（对应 $x_{\min}$ 抬到 8 以上），
Hill 估计的漂移就会超过零点九、任何单点 $\hat\alpha$ 都不该被引用来作幂律证据。
本篇因此给出一条可复算的判据而不是一个形容词：
\textbf{尾样本量在百条量级及以下、且 $\hat\alpha$ 随 xmin 漂移接近或超过一个单位时，
停止使用"无标度"这一标签。} 这是一条负结果，但它可迁移——
换到另一部叙事、另一套抽取，只要尾样本落进同一量级，收手的条件同样成立。
\item \textbf{L5 类别标签是粗标签，卡方只报一次。}
关系与事件各有 12 与 11 个类别，其中总有占两三成的粗类（「其他」、未分类）几乎不动，
没有区分力，本篇把它们留在表里当分母与核对项。
卡方独立性检验在本篇只做一次（不做多重比较后的挑选），
因为三期三刀是先验固定的，不是从数据里搜出来的；
因此 p 值不需要作多重校正，但也不该被当成"类别迁移被证实"——
它只否定了独立，没肯定强度，强度由 V=0.259 说了。
\item \textbf{L6 时期切分是先验的，占比迁移的显著性依赖这一刀。}
P1/P2/P3 的边界取自王朝分期，不来自聚类；换一套切点（尤其是 1600 这一刀），
军事战争与政治斗争的走向未必完全保持。本篇因此报的是"在这三刀下的位移"，
不声称发现了一个断代点。
\item \textbf{L7 自审不是跨模型对抗评审。}
本机无第二模型端点，评审由同一模型多轮盲读完成；
已知后果是写作轮曾把未核验的文献键写进正文、盲读轮才拦下。
自审能抓内部不一致，抓不到共同先验盲区。
\end{itemize}

需要强调的是：以上限制没有一条是再跑一遍数据就能消掉的。
L1、L2 是同源投影的固有位置，L4 是样本量的固有位置，L6 是先验切分的固有位置。
本篇的处理是把它摊开并给出可迁移的判据——能降到差别的，就不叫效应；
不能降的，就明确写为边界。
